\chapter{从现实问题到机器学习问题：建模的艺术}

AI 最大的瓶颈往往不在算力，而在于使用者能否把问题说清楚。前两章解决了“如何向 AI 提问”的问题：理解工具能力边界、掌握 Prompt 的构成要素和写作技术。本章进入更上游的位置——在动笔写第一条 Prompt 之前，使用者必须先回答一个更根本的问题：你真正想解决的是什么问题？

这个问题比听起来难得多。现实中的任务往往是模糊的、充满矛盾的，甚至连利益相关方自己也说不清楚想要什么。“降低城区拥堵”可以被翻译成几十种不同的机器学习任务，每种都有不同的数据要求、评价标准和错误成本。选错建模方向，再精巧的 Prompt 技术也无济于事。

本章讨论如何将模糊的现实问题转化为可计算的机器学习问题，并将这些建模决策通过 Prompt 传递给 AI。这一过程——问题建模——是整个协作链条中最依赖人类判断、AI 最难替代的环节。

问题建模的核心不是“怎样套用一个模型”，而是建立从决策目标到可观测数据、再到评价标准的完整链条。数据科学方法论通常把业务理解与问题定义置于算法选择之前，因为一个预测结果只有在能改变决策时才产生价值\cite{provost2013data}。这也意味着，项目开始时应当允许得出“当前不需要机器学习”或“现有数据不足以回答问题”的结论。

% ----------------------------------------------------------------
\section{为什么问题建模是人类不可替代的核心能力}
% ----------------------------------------------------------------

\subsection{现实问题天然是模糊的}

工程和研究中的真实任务，从来不以“机器学习问题”的形式出现。它们往往以如下方式被提出：“我们希望缓解早高峰的校园周边拥堵”；“能不能用数据分析一下，到底是哪些原因导致事故频发”；“我们想提高公交系统的服务水平”。这些陈述有几个共同特点：目标定义含糊、缺乏可量化的成功标准、利益相关方的诉求相互交叠，以及错误的代价因立场而异。

\textbf{目标含糊}是最直接的问题。“缓解拥堵”至少可以被解读为：减少拥堵的发生频率、缩短拥堵的持续时长、降低平均延误、减少极端延误事件，或者在有限道路资源下使更多人准时到达目的地。每一种解读对应不同的建模目标，也需要不同的数据和评价指标。

\textbf{标签不明确}是第二层困难。机器学习模型需要一个可以量化的学习目标，而真实任务中这个目标往往不是现成的。“服务水平”不是一个字段，它是由候车时间、车内拥挤程度、准点率和班次间隔共同构成的综合评价。在决定如何建模之前，必须先决定如何测量。

\textbf{利益相关方多}带来了冲突的目标。同一个交通系统中，乘客希望最短等待时间，运营商希望最低运营成本，监管方希望覆盖更多区域，周边居民希望减少噪音和拥堵外溢。没有一个模型能同时最优化所有人的目标。选择建模目标，本质上是在做一个价值取舍的决定，而这类决定不能外包给 AI。

\textbf{错误成本因人而异}是第四层障碍。事故预测中，漏报一次严重事故的代价与误报十次轻微风险的代价，对于安全工程师和道路运营商来说有完全不同的权重。这些权重不写在数据里，也不会从模型中自动浮现，只有决策者才能给出。

\subsection{AI的能力边界：能执行定义好的问题}

AI 在建模任务中的能力边界，可以用一句话概括：\textbf{AI 能够高效执行已被清晰定义的问题，但无法替使用者决定什么问题值得解决，以及什么结果才算正确。}

给定一个清晰的建模定义，AI 能够选择合适的模型结构并生成训练代码，处理数据、构造特征、划分训练集和验证集，计算指定评价指标并生成对比报告，比较多种方法的性能并解释差异，并根据结果提出下一步探索方向。但在给出清晰定义之前，AI 会填补假设——它会选择最常见的任务类型、最常用的评价指标、最普遍的数据划分方式，而这些默认选择在特定领域下往往是错误的。第二章已经说明：当 Prompt 缺失关键信息时，模型会用训练数据中最常见的模式填补空白。建模层面的问题在于，这些“最常见模式”是通用场景下的默认值，而交通运输工程有自己的规则、定义和约束。

此外，预测准确并不自动意味着解释正确。Breiman 对统计建模两种文化的讨论指出，强调数据生成机制的解释性建模与强调泛化性能的算法建模关注不同问题\cite{breiman2001statistical}。在实际项目中，两者都可能有价值，但不能互相冒充：用于预测拥堵的高性能模型，不一定能够回答修建新道路是否会缓解拥堵；发现事故与天气相关，也不等于改变天气变量即可估计政策干预效果。

\subsection{同一份数据，十几种不同的ML任务}

以“交通事故数据”为例，直观说明建模决策的核心作用。假设拥有一份包含事故时间、地点、类型、严重程度、道路条件和气象信息的历史记录数据集，同一份数据可以被翻译成以下截然不同的机器学习任务：

\begin{enumerate}
  \item \textbf{路段事故率预测（回归）：}预测未来一个月内，每条路段发生事故的概率，为巡查资源分配提供依据。
  \item \textbf{事故严重程度分类（多分类）：}给定事故发生后的初步报告，自动判断严重级别（轻微/一般/重大），用于快速触发不同级别的应急响应。
  \item \textbf{高风险时段检测（二分类）：}判断某路段在特定时段是否属于高风险状态，辅助动态限速决策。
  \item \textbf{事故黑点聚类（无监督聚类）：}发现事故频发的空间集群，以便确定改造优先级。
  \item \textbf{事故影响因素分析（可解释建模）：}识别哪些路段设计特征或气象条件与事故率显著相关，为工程改造提供依据。
  \item \textbf{干预效果评估（因果推断）：}评估某次道路改造措施在实施前后是否真正降低了事故率，排除季节性和外部变化的干扰。
  \item \textbf{事故漏报异常检测（异常检测）：}识别数据中异常低的事故记录密度，判断是否存在漏报或数据采集问题。
  \item \textbf{行程安全评分（排序学习）：}对同一起终点的多条路径按安全性排序，为驾驶导航提供参考。
\end{enumerate}

这八种任务都在同一份数据上进行，却需要完全不同的目标定义、标签设计、模型结构和评价方式。没有一种是“自然正确”的选择——选哪一种，取决于真实业务目标、可用数据的质量和范围、结果将如何被使用，以及什么样的错误是不可接受的。这个选择必须由任务负责人做出，不能外包给 AI。

% ----------------------------------------------------------------
\section{现实问题到机器学习问题的翻译框架}
% ----------------------------------------------------------------

将现实问题翻译为机器学习问题，需要在五个维度上做出明确决策。这五个维度共同决定了模型的形态、数据需求、评价方式和可用性边界。每个维度的决策都是价值判断，不是技术判断。

\subsection{第一维度：预测对象是谁}

机器学习模型的输出针对的是一个特定的\textbf{分析单元}。在交通运输领域，不同的分析单元对应着截然不同的建模逻辑。\textbf{路段（Link）}由检测器数据或路网划分定义，适合流量预测、速度预测、拥堵状态分类等任务。\textbf{交叉口（Intersection）}适合信号配时优化、冲突事件分析、行人过街风险评估。\textbf{起终点对（OD Pair）}适合出行需求预测、行程时间估计、路径选择建模。\textbf{车辆（Vehicle）}适合驾驶行为识别、异常行为检测、公交运营监测。\textbf{乘客（Passenger）}适合乘车体验分析、换乘行为建模、票务异常检测。\textbf{区域（Zone）}适合出行需求热点分析、停车需求预测、共享出行调度。\textbf{道路网络（Network）}则适合全网流量状态估计、事件传播分析、韧性评估。

分析单元的选择直接影响数据聚合方式和可用模型。以“公交需求预测”为例：如果分析单元是站点，任务是预测每个站点的上客人数；如果分析单元是 OD 对，任务是预测每条出行路径上的需求量；如果分析单元是线路，任务是预测线路总客流。这三种定义需要不同的数据基础，也回答不同的业务问题。分析单元的选择还必须与决策权力边界匹配：如果政策干预的最小执行单元是区，则以路段为粒度的精细预测结果无法被直接使用。

\subsection{第二维度：任务类型是什么}

确定了分析单元之后，需要决定对该单元执行什么类型的计算任务。交通运输领域最常见的任务类型，各有其适用范围与逻辑。\textbf{预测（Prediction/Forecasting）}基于历史数据预测未来某个时点或时段的指标值，例如预测路段未来 15 分钟的平均速度，或预测节假日期间高速公路的小时流量。\textbf{分类（Classification）}将输入划分为若干预先定义的类别，例如判断驾驶行为是否危险，或对路段状态进行拥堵等级标注。\textbf{聚类（Clustering）}在无标签的情况下发现数据中的自然分组结构，例如按出行规律对乘客进行分群，或按运行特征对公交站点进行分组。\textbf{异常检测（Anomaly Detection）}识别与正常模式明显偏离的观测，例如检测检测器数据中的设备故障，或发现流量数据中的突发事件。\textbf{因果推断（Causal Inference）}估计某种干预或政策的真实效果，与单纯的相关关系分析有本质区别，例如评估拓宽路肩是否降低了事故率，或分析收费政策变化是否改变了出行需求。\textbf{优化（Optimization）}在约束条件下寻找最优决策方案，例如在给定班次数量和车辆容量条件下最优化公交发车间隔。\textbf{生成与仿真（Generation/Simulation）}生成符合真实分布特征的合成数据，或模拟复杂系统在不同条件下的行为，例如生成用于测试自动驾驶系统的边界场景。

任务类型的选择应当由业务目标驱动，而不是由数据的便利性或技术的流行程度驱动。一个常见的错误是，因为手头有大量时序数据，便自然地滑向预测任务，而实际业务需要的是识别异常或评估政策效果。任务类型一旦选错，再精良的模型也无法提供有用的答案。

\subsection{第三维度：标签从哪里来}

对于监督学习任务，标签是整个建模过程中最关键、最容易被轻视的环节。标签定义了模型的学习目标，标签的质量直接决定了模型能学到什么。

\paragraph{谁定义标签}

标签的定义是一个价值判断，不是技术判断。以事故严重程度为例，同一次事故可能得到不同标签。按医疗损失定义，涉及人员住院治疗的为严重，仅需急救处理的为一般；按道路影响定义，导致车道封闭超过 2 小时的为严重，不封闭车道的为一般；按财产损失定义，损失超过特定金额的为严重；按法规定义，则依照当地交通法规中的事故分级标准。这四种定义可能对同一次事故给出不同的标签。使用哪种定义，不是数据科学问题，而是领域和政策问题，必须由领域专家或决策者明确后写入 Prompt。

\paragraph{怎么获取标签}

标签的获取方式决定了标签的可靠性和规模。\textbf{官方记录}（警方事故记录、医院急救数据、保险理赔记录）通常质量较高，但存在漏报和记录延迟问题，且覆盖的事故级别可能有下限，轻微事故通常不报警。\textbf{传感器触发}（检测器速度骤降、摄像头异常事件）覆盖面广，但判别规则简单，噪声较多，需要人工核查。\textbf{人工标注}由领域专家对历史记录进行标注，标签质量高，但成本高、规模有限，且不同标注者之间可能存在分歧。\textbf{众包标注}由网约车司机、志愿者或乘客提供报告，规模大，但质量参差不齐，存在系统性偏差——热门路段被过度报告，偏远路段被忽视。

\paragraph{标签质量如何影响模型}

标签错误不是随机噪声，往往具有系统性规律。如果某类事件天然更难被记录（如夜间轻微碰擦），则这类事件的标签缺失不是随机缺失，而是与时间、地点、严重程度相关的系统性缺失。在这种情况下，用有偏标签训练的模型会系统性地低估夜间轻微事故的风险，即使训练集指标看起来不错。

在向 AI 提交建模任务之前，必须能够回答：标签由谁定义？采用什么定义标准？数据集中标签的完整性如何？是否存在系统性缺失或错误？

\subsection{第四维度：时间与空间粒度}

粒度选择直接影响模型能够回答的问题，也影响数据需求和计算成本。粒度选择错误，会导致模型在技术上运行正常，但无法支持实际决策。

\paragraph{时间粒度}

时间粒度的选择决定了模型服务于哪个层次的决策。\textbf{5分钟级}适合实时信号控制、短期行程时间估计、突发事件检测，但数据量大，对传输和计算的要求高，模型需要处理短时噪声。\textbf{小时级}适合日内需求预测、停车管理、出租车和网约车调度，在大多数应用中是时序分析的基础粒度。\textbf{日级}适合长期流量趋势分析、政策效果评估、年度运营规划，通常与气候、假期、活动等宏观因素关联更强。\textbf{周级/月级}适合中长期规划、基础设施需求分析、季节性规律研究，在这一粒度上随机波动被平滑，长期趋势更为清晰。

以公交客流预测为例：5 分钟粒度的预测结果可以直接触发动态发车决策；小时粒度的结果更适合前一天制定发班计划；日粒度的结果服务于月度运力配置。三种粒度服务于不同的决策时间窗口，不能互换。

\paragraph{空间粒度}

空间粒度同样需要与决策需求匹配。\textbf{点位（传感器/站点）}是最细的空间粒度，对应单个检测器或公交站点，适合需要精确位置信息的任务。\textbf{路段}将点位数据聚合到路段层面，是大多数路网分析的基本单元。\textbf{交通分析小区（TAZ）}是城市交通规划中的标准空间划分，适合出行需求建模。\textbf{行政区划（街道/区县）}适合政策分析和行政管理决策支持。\textbf{格网（Grid）}将城市划分为均匀网格，适合需要空间一致性的深度学习方法，如卷积神经网络。

\subsection{第五维度：错误成本的非对称性}

这是五个维度中最容易被技术导向的使用者忽视、但往往对实际决策影响最大的一个维度。

大多数机器学习课程使用对称的评价指标——均方误差、准确率、F1 分数。这些指标隐含了一个假设：不同类型的错误代价相同。在交通运输工程中，这个假设几乎总是错误的。以\textbf{桥梁结构安全预测}为例，漏报一次真实的安全隐患（将危险桥梁预测为安全）可能导致灾难性事故，而误报一次安全隐患（将正常桥梁预测为危险）只会造成不必要的检查费用；在安全领域，漏报的代价通常远高于误报，因此应赋予召回率更高的权重。在\textbf{公交车辆调度}场景中，预测客流不足导致配车过多浪费运营成本，而预测客流过高导致配车不足则让乘客无法乘车，在面向乘客的服务中，漏报高峰通常造成更严重的服务下降。\textbf{自动测速执法}则面临另一种不对称：误判正常车辆超速会损害公信力、引发投诉，漏判真实超速车辆则减少违法行为的威慑效果，两类错误的政治和法律成本截然不同，需要在模型设计时明确权衡。

错误成本的非对称性必须在建模阶段解决，而不是在模型训练完成后再补救。解决方式包括：

\begin{enumerate}
  \item \textbf{选择合适的评价指标：}使用精确率（Precision）和召回率（Recall）代替准确率；在回归任务中使用非对称损失函数；设计反映真实业务代价的自定义指标。
  \item \textbf{调整分类阈值：}在预测概率已经可靠的前提下，通过调整决策阈值改变精确率和召回率之间的权衡，使其与业务优先级一致。
  \item \textbf{使用非对称损失函数训练：}在模型训练时直接引入不同类型错误的代价权重，而不是事后调整。
  \item \textbf{将成本量化并报告：}在评估报告中明确展示不同错误类型的数量及其估计代价，使决策者能够在真实成本框架下判断模型是否可用。
\end{enumerate}

% ----------------------------------------------------------------
\section{机器学习问题的主要范式}
% ----------------------------------------------------------------

在确定了分析单元、任务类型、标签、粒度和错误成本之后，需要选择具体的机器学习范式。不同范式的适用条件、数据要求和局限性差异显著。本节介绍交通运输领域最常用的几类范式，并特别说明交通数据的结构特殊性。

\subsection{监督学习}

监督学习是交通运输领域应用最广泛的范式。其核心特点是：存在明确的预测目标（标签），模型通过学习输入特征与标签之间的映射关系来生成预测。

\paragraph{回归任务}

回归任务预测连续数值。在交通运输领域，典型应用包括：\textbf{交通流量预测}，预测路段未来某时段的车辆通过数量，服务于信号配时和动态路径诱导；\textbf{行程时间估计}，预测从起点到终点的行驶时长，服务于导航和物流配送计划；\textbf{停车需求预测}，预测特定时段内停车场的占用率，服务于停车诱导系统；以及\textbf{公交到站时间预测}，预测公交车到达各站点的时间，服务于乘客信息系统。

回归任务的评价通常使用均方误差（MSE）、平均绝对误差（MAE）或平均绝对百分比误差（MAPE）。选择哪种指标，应根据错误成本结构决定：MAE 对大误差和小误差等权处理；MSE 对大误差施加更高惩罚，适合对极端误差特别敏感的场景；MAPE 将误差标准化为百分比，适合需要跨路段比较相对误差的情形。

\paragraph{分类任务}

分类任务将输入划分为预先定义的类别。典型应用包括：\textbf{事故等级判断}，给定事故描述和现场信息，预测事故属于哪个严重级别，触发对应的应急资源调配；\textbf{拥堵状态标注}，将路段当前状态分为畅通、缓行、拥堵，服务于路况信息发布；\textbf{驾驶行为识别}，基于车载传感器数据，将驾驶行为分类为正常、激进、疲劳等模式，服务于安全驾驶辅助；以及\textbf{公交准点判断}，预测某次班次在当前运行条件下是否会准时到达终点，支持调度员早期干预。

分类任务的评价需要区分不同错误类型的代价。准确率（Accuracy）在类别不平衡时会产生误导——如果 95\% 的情况都是“无事故”，一个永远预测“无事故”的模型准确率为 95\%，但毫无实用价值。应优先使用精确率、召回率、F1 分数，以及混淆矩阵和 ROC 曲线。

\subsection{无监督学习}

无监督学习不依赖人工标注的标签，而是直接从数据的内在结构中发现规律。这在标签成本高昂或根本无法定义标签的场景中特别有价值。

\paragraph{聚类分析}

聚类的目标是将相似的对象归组，发现数据中自然存在的分组结构。在交通运输领域，聚类有几类典型用途：\textbf{出行模式挖掘}，基于乘客的出行时间、频率和路径特征将乘客聚类为不同出行模式群体，如工作日通勤者、偶发出行者、夜间出行者，支持差异化的票价和服务设计；\textbf{道路功能分类}，基于时序流量特征将路段聚类为功能相似的类型，如商业区集散路段、通勤走廊、工业区货运通道，支持精细化管理；以及\textbf{事故多发区域识别}，将地理空间上的事故记录聚类，识别需要重点干预的空间热点，辅助改造优先级排序。

聚类结果没有唯一“正确”答案，质量评估需要结合领域知识。一个在统计指标上最优的聚类结果，可能在业务上毫无意义。因此，聚类分析需要领域专家参与对结果的解读和验证。

\paragraph{异常检测}

异常检测识别与正常模式显著偏离的观测，适用于没有完整事故标签但需要实时发现问题的场景。常见应用包括：\textbf{传感器故障检测}，在道路检测器数据中识别连续零值、恒定值或统计特征突变等异常模式，区分设备故障和真实交通状态变化；\textbf{交通事件实时检测}，在视频或流量数据中识别速度突然下降、流量骤变等异常，触发事件确认流程；以及\textbf{GPS 轨迹异常识别}，在公交 GPS 数据中识别偏离既定路线、长时间停顿或行驶速度异常的班次，支持运营监控。

异常检测的一个关键挑战是定义“正常”。交通数据有强烈的时间周期性（工作日与周末不同、早晚高峰不同、季节不同），如果不考虑时间背景，一个平日高峰期的正常流量值，在夜间可能被误标记为异常。

\paragraph{降维与可视化}

降维方法将高维数据压缩到低维空间，用于探索、可视化和特征提取。在交通运输研究中，常见用途包括：将多个检测器的流量时序数据压缩为少数几个主成分，揭示全网流量的主要变化模式；将高维驾驶行为特征压缩后可视化，辅助人工理解不同驾驶风格的分布。

\subsection{强化学习}

强化学习（Reinforcement Learning, RL）通过智能体与环境的持续交互来学习决策策略，适用于序列决策问题，即当前的决策会影响未来状态的场景。在交通运输领域，RL 有几类典型的应用场景。\textbf{信号配时优化}中，智能体控制路口信号灯，根据当前排队状态和流量选择配时方案，以最小化全网延误为目标，其优势在于能够适应实时变化的流量，而不依赖预设的固定方案。\textbf{动态定价}场景中，智能体在拥堵收费或网约车定价时根据当前供需状态调整价格，以在服务水平和收益之间实现动态平衡。\textbf{自动驾驶决策}则在微观驾驶层面让智能体学习换道、超车、跟车等操作的决策策略，在复杂交通环境中保持安全和效率。

强化学习的部署复杂性远高于监督学习。主要挑战包括：在真实道路上试错的代价极高，通常需要先在仿真环境中训练，再谨慎地迁移到真实系统；奖励函数的设计往往比模型本身更难，设计不当的奖励会导致模型学到满足形式目标但违背真实目标的策略；以及训练和收敛所需的样本量通常非常大。对于大多数交通运输工程项目，强化学习不是首选方案。在决定使用 RL 之前，应先确认监督学习能否以更低的代价解决问题。

\subsection{交通数据的特殊结构}

交通数据不是独立同分布的样本集合，而是具有复杂时空结构的序列。正确处理这种结构，是交通数据建模区别于通用机器学习教程的核心所在。

\paragraph{时序结构}

交通数据具有两种典型的时序规律。一是\textbf{周期性（Periodicity）}：工作日早晚高峰、周末出行模式、节假日特殊规律在特征工程中必须显式建模，否则模型无法区分“工作日早高峰的正常拥堵”和“同时刻的异常事件”。二是\textbf{趋势性（Trend）}：城市发展、道路改扩建、土地利用变化会导致长期趋势变化，在长时间跨度的数据上训练的模型需要考虑早期数据是否仍然代表当前系统状态。

时序数据最严重的建模错误是\textbf{数据泄露}：将未来才能获得的信息用于预测当前时刻。例如，用明天的气象数据预测今天的事故风险，在测试集上看起来效果很好，在真实部署中完全失效。预防数据泄露的核心原则是：特征必须只来自预测时刻之前已经可以获得的数据，验证集必须严格按时间顺序位于训练集之后。

\paragraph{空间依赖}

交通现象具有强烈的空间依赖性。一条路段的拥堵会向上游传播；一个路口的信号变化会影响相邻路段的排队。这种空间相关性意味着：将每个路段视为独立样本进行建模，会忽略相邻路段之间的因果关系；空间上相邻的测试样本之间可能存在信息泄露，导致交叉验证结果偏乐观；对单一路段建立的模型，在相邻路段上的泛化性能可能很差。

\paragraph{图结构}

道路网络和公交线网具有天然的图结构：节点是路口或站点，边是路段或线路。这种图结构对于需要建模网络传播效应的任务（如全网流量估计、事件影响范围预测）非常重要。图神经网络（GNN）是近年来专为处理此类结构而发展的方法，但在引入复杂模型之前，应首先评估图结构是否对当前任务真正必要。

\subsection{在同一项目中组合使用多种范式}

真实项目往往不是单一范式的简单应用，而是多种范式的组合。\textbf{异常检测与监督分类}的组合是一种常见模式：先用异常检测方法识别可疑样本，再由人工确认标签，再用监督分类学习事件类型，从而在标签资源有限的情况下，将人工标注集中在最有价值的样本上。\textbf{聚类与监督回归}的组合则先用聚类将路段按功能分组，再为每类路段分别训练独立的预测模型，利用聚类发现的分组结构减少单一模型需要处理的异质性。\textbf{预测与优化}的组合先用监督学习预测未来流量，再将预测结果作为优化问题的输入，求解最优信号配时或调度方案，预测模型和优化模型分别独立建立，接口处通过预测结果连接。

组合使用多种范式时，需要特别注意各阶段之间的误差传播：第一阶段的误差会被第二阶段放大。在设计系统时，应为每个阶段独立设计评价标准，而不是只关注最终输出的表现。

% ----------------------------------------------------------------
\section{问题定义的常见错误}
% ----------------------------------------------------------------

以下四类错误在实际项目中反复出现，且往往在投入大量建模工作之后才被发现。提前识别这些错误，可以在最低成本时纠正方向。

\subsection{以“准确率最高”为唯一目标}

\textbf{错误表现：}在确定建模目标时，只写下“建立一个尽可能准确的预测模型”，而不考虑哪种错误更不可接受。

\textbf{后果：}模型优化过程向主流样本倾斜，少见但高影响的情况被忽视。以事故严重程度预测为例，如果严重事故只占样本的 5\%，最大化总体准确率的模型会倾向于将所有样本预测为“一般”或“轻微”，在不严重漏掉严重事故的情况下实现高准确率，但这正是最危险的结果。

\textbf{修复方向：}在建模目标中明确说明不同类型错误的相对权重，选择与错误成本结构一致的评价指标，并在模型设计时通过调整阈值或损失函数体现这种权重。

\subsection{目标代理：优化了代理指标，偏离了真实目标}

\textbf{错误表现：}真实目标无法直接量化，于是选择一个与真实目标相关的代理指标，并在整个建模过程中将其作为唯一优化目标。最终，模型在代理指标上表现优秀，但真实目标并未改善。

\textbf{交通领域例子：}某市希望“提升公交出行体验”，由于无法直接量化乘客体验，选择“准点率”作为代理指标并集中优化。模型成功将准点率从 72\% 提升到 89\%，但乘客满意度调查没有改善。原因是：为了提高准点率，调度策略倾向于取消晚点班次，而不是通过真正改善运营效率来保证准点，这在统计上提高了准点率，但在乘客体验上适得其反。

\textbf{修复方向：}在确认代理指标时，明确说明代理指标与真实目标之间的关系假设，以及这个假设在什么条件下可能失效。在项目过程中，定期检查代理指标的改善是否真正伴随了真实目标的改善。

\subsection{把相关关系当因果：发现规律后直接采取政策干预}

因果问题必须明确干预、对照和目标人群，并讨论混杂因素、选择偏差及识别假设。因果推断教材反复强调，“观察到什么”与“如果采取另一项行动会发生什么”是两类不同的问题\cite{hernan2020causal}。让 AI 自动选择一个预测模型再解释特征重要性，不能填补这两类问题之间的逻辑缺口。

\textbf{错误表现：}在数据中发现了两个变量之间的强相关关系，直接将其解读为“A 导致了 B”，并据此设计政策干预，期待通过改变 A 来影响 B。

\textbf{交通领域例子：}分析发现某区域共享单车停靠密度与该区域步行街的客流量高度正相关。据此，决策者在另一个区域大量投放共享单车，期望带动客流。实际上，共享单车和客流可能共同受到该区域商业发展程度的影响——商业活跃的区域既吸引出行（带来更多单车需求），也带来更多步行客流。单车投放并不是客流的原因，而是与客流同时出现的结果。

要将相关分析升级为因果推断，需要使用专门的方法：随机对照实验（最可靠，但成本高）、双重差分法（Difference-in-Differences）、断点回归（Regression Discontinuity）或工具变量法（Instrumental Variables）。这些方法各有适用条件，选择时需要认真考虑数据结构和识别策略。

\subsection{把技术可行性当业务可行性}

\textbf{错误表现：}模型在历史数据上训练完成，测试集指标令人满意，但结果无法被实际业务流程所使用。

\textbf{交通领域例子：}这类情况以多种面貌出现。交通事故预测模型达到了 85\% 的准确率，但部署时发现预测所需的实时气象数据接口不稳定，每天有 3—4 小时无法获取，导致系统在高风险时段反而不可用。公交调度优化模型在仿真中将总延误降低了 18\%，但实施后发现优化方案要求驾驶员班次调整频率超过了劳动合同规定，无法推行。路网流量估计模型在训练期数据上表现优秀，但实际使用时发现，模型所依赖的部分检测器在恶劣天气下故障率高达 40\%，在最需要预测的时段恰好数据最不完整。

\textbf{修复方向：}在确定建模方向时，就应当与运营方讨论系统边界：预测结果需要多快？数据接口的可靠性如何？预测结果将以什么形式被人使用或触发自动决策？这些限制条件应当在建模设计阶段明确，而不是在系统开发完成后再发现。

% ----------------------------------------------------------------
\section{将建模决策转化为Prompt}
% ----------------------------------------------------------------

前几节建立了问题建模的分析框架。本节回到操作层面，说明如何将这些建模决策转化为可传递给 AI 的 Prompt，使 AI 能够在正确的问题定义下工作。

\subsection{把五个维度写进Prompt}

五个维度的建模决策都应当在 Prompt 中显式声明。分析单元的声明方式例如：“本任务的分析单元是路段，每条路段由检测器编号标识，数据已按路段聚合，不需要从传感器点位数据重新聚合。”任务类型的声明方式例如：“任务类型是时序预测（回归）：给定路段过去 60 分钟的流量序列，预测接下来 15 分钟的平均流量。不需要做分类或聚类，不需要识别异常，只需要生成点预测值。”标签定义的声明方式例如：“预测目标字段是 \texttt{flow\_15min}，单位为辆/15分钟，由检测器计数直接汇总得到，已通过质检，不存在系统性偏差。标签是连续数值，不需要离散化。”粒度的声明方式例如：“时间粒度为 15 分钟，空间粒度为单条路段。训练数据覆盖 2022 年全年，测试数据为 2023 年 1—3 月，两者严格按时间顺序划分，不允许随机划分。”错误成本的声明方式例如：“本项目对高流量时段的预测精度要求更高，因为预测误差在高峰期对信号配时的影响更大。评价指标使用加权 MAPE，其中早晚高峰时段（7:00—9:00 和 17:00—19:00）的权重为非高峰时段的 3 倍。”

\subsection{把领域隐含知识传递给AI}

领域专家通常掌握大量 AI 无法自动知道的隐含知识。这类知识至少涵盖五个方面：数据采集机制（检测器的工作原理、常见故障模式、数据缺失的规律和含义）；业务约束（运营规则、劳动合同限制、法规要求、决策权边界）；领域定义（行业内对“拥堵”、“准点”、“服务水平”等概念的标准定义或组织内部定义）；数据历史（数据集覆盖期间是否发生过重大事件——道路施工、政策变化、节假日调整——以及这些事件是否应当被排除或特殊处理）；以及先验知识（通过领域经验知道某些变量应当具有的方向关系或数量级，可以作为模型合理性检验的参考）。

应在 Prompt 的“约束”或“上下文”部分明确声明这些知识，而不是期待 AI 自行推断。以下是一个领域背景说明的示例：

\begin{quote}
\textbf{【领域背景说明】}

本数据来自感应线圈检测器，每 5 分钟汇总一次计数。以下是使用本数据必须了解的领域知识：

（1）检测器在雨天的故障率约为晴天的 2 倍，表现为 \texttt{speed} 字段连续出现 0 值，而 \texttt{flow} 字段同时归零。两者同时为零的连续记录（超过 3 步）应视为设备故障，而非真实的零流量状态。

（2）本路段每周二上午 10:00—12:00 有固定的货车限行规定，该时段流量会显著偏低，不代表拥堵缓解，也不代表需求下降。特征工程时应构造“限行时段”标志位。

（3）数据集涵盖 2021 年 7—8 月的疫情管控期，该期间数据显著异常，建议在训练集中剔除，或至少单独评估其影响。

（4）“拥堵”在本项目中定义为：平均速度低于 20 km/h 且持续超过 15 分钟。这与部分文献中使用的 30 km/h 阈值不同，请不要使用文献中的通用定义。
\end{quote}

\subsection{完整的建模问题定义Prompt示例}

以下是一个将上述所有要素整合的完整 Prompt 示例，场景是预测城市快速路某路段的短期流量：

\begin{quote}
\textbf{角色：}你是一位交通数据分析工程师，正在为城市交通管理中心开发实时流量预测系统的核心模型。

\textbf{任务背景：}我需要建立一个短期流量预测模型，预测结果将每 15 分钟更新一次，并被实时信号控制系统使用。模型推理时间不超过 1 秒。

\textbf{数据说明：}
\begin{itemize}
  \item 数据文件：\texttt{detector\_flow\_2023.csv}，共 20 个路段，时间范围 2023 年 1—10 月，时间粒度 15 分钟。
  \item 字段：\texttt{timestamp}（格式 YYYY-MM-DD HH:MM）、\texttt{detector\_id}（路段编号，字符串）、\texttt{speed}（km/h，浮点数，-1 表示设备离线）、\texttt{flow}（辆/15分钟，整数）、\texttt{occupancy}（占有率，0—100 的浮点数）。
  \item 已知质量问题：\texttt{speed} 为 -1 的记录是设备离线状态，\texttt{flow} 在此时通常为 0 但不可靠，均应标记为缺失值。速度超过 130 km/h 的记录属于测量错误，应删除。
\end{itemize}

\textbf{建模定义：}
\begin{itemize}
  \item 预测目标：给定当前时刻及过去 4 个时间步（即过去 1 小时）的 \texttt{flow} 值，预测下一个时间步（未来 15 分钟）的 \texttt{flow} 值。
  \item 训练集：2023 年 1—8 月；验证集：9 月；测试集：10 月。三个集合严格按时间顺序划分，禁止随机划分。
  \item 所有特征只能使用当前时刻之前已经可以获得的数据，禁止使用未来信息。
\end{itemize}

\textbf{领域约束：}
\begin{itemize}
  \item 早晚高峰（7:00—9:00 和 17:00—19:00）的预测精度更重要，这两个时段对信号控制的影响最大。
  \item 本路段每周三 22:00—次日 6:00 有夜间施工，该时段数据不代表正常运营状态，可在特征工程中构造标志位，或在评估时单独处理。
\end{itemize}

\textbf{评价指标：}使用加权 MAPE 作为主要指标，其中高峰时段权重为 3，非高峰时段权重为 1。同时报告 MAE，用于直观理解误差大小。

\textbf{任务：}请按以下步骤进行，每步完成后输出验证信息，确认无误后再继续：

\begin{enumerate}
  \item 读取数据，输出每个字段的非空率、基本统计（最小值、最大值、均值）和时间范围，检查是否与预期一致。
  \item 按上述约束处理缺失值和异常值，输出处理前后每路段的数据量和缺失率变化。
  \item 构造时间特征（小时、星期几、是否高峰时段、是否施工时段）和滞后特征（过去 4 步的 \texttt{flow} 值），输出特征列表和样例行，确认没有未来信息。
  \item 按时间顺序划分训练集、验证集和测试集，输出各集合的起止时间和样本量。
  \item 训练一个基准模型（历史同时刻均值基准），计算验证集上的加权 MAPE 和 MAE。
  \item 在基准之上，尝试一个线性模型（岭回归）和一个集成模型（随机森林），比较三种方法在验证集上的加权 MAPE。所有模型使用相同的训练集和验证集，结果放入同一张对比表。
\end{enumerate}

暂时不要替我决定最终使用哪个模型；等输出对比结果后我自行判断。
\end{quote}

这个 Prompt 将五个维度的建模决策全部显式化，将领域隐含知识转化为明确规则，并要求 AI 分步执行、逐步验证。它比一般 Prompt 更长，但这种长度换来的是大幅减少后续的反复调试，以及更清晰的人机分工：AI 负责执行每一步，人负责在每步之间确认结果是否符合预期。

% ----------------------------------------------------------------
\section{贯穿案例}
% ----------------------------------------------------------------

\begin{quote}
\textbf{【案例占位】}

本节将通过一个完整的交通运输工程案例，展示从模糊现实问题到可执行建模定义的全过程。案例将覆盖：任务背景与利益相关方分析、五个维度的逐一决策过程、建模过程中遭遇的典型问题（标签获取困难、错误成本争议、空间粒度选择），以及最终形成的完整建模问题定义 Prompt。

\textbf{【待确定内容】}案例的具体背景（候选方向：高速公路事故快速检测、城市公交高峰期调度优化、地铁客流预测与运力配置）；案例涉及的真实数据结构；五个维度各自的决策过程和关键分歧点；最终形成的完整建模定义文档。
\end{quote}

% ----------------------------------------------------------------
\section{实战练习}
% ----------------------------------------------------------------

\subsection*{练习一：识别任务类型}

以下是五个来自交通运输工程实践的任务描述。对每个描述，完成以下判断：（1）该任务属于哪种机器学习任务类型（预测、分类、聚类、异常检测、因果推断、优化）；（2）分析单元是什么；（3）可能的标签来源是什么，主要的质量风险是什么；（4）最重要的错误成本不对称体现在哪里。

\smallskip
\textbf{任务一：}某高速公路管理中心希望在事故发生后 5 分钟内自动识别事故位置，触发应急响应，避免二次事故。现有数据包含沿线检测器的流量和速度记录，以及历史事故报告。

\smallskip
\textbf{任务二：}某城市公共交通公司希望了解，将公交卡优惠政策（早高峰 8 折）在试点线路上实施三个月后，是否真正增加了年轻通勤族的公交出行频率。

\smallskip
\textbf{任务三：}某共享单车运营商希望将城市内的停车区域按使用模式分群，以便为不同类型区域制定差异化的车辆调配和维护策略。

\smallskip
\textbf{任务四：}某路段的交通工程师发现某个检测器近期报告的数据与周边检测器明显不一致，怀疑设备存在故障或外部干扰，需要自动化识别此类问题。

\smallskip
\textbf{任务五：}某城市希望在已有公交运营数据的基础上，预测明天各主要换乘站点早高峰期间的客流量，以便提前部署安保和服务人员。

\subsection*{练习二：发现建模定义中的问题}

以下是一个已经写好的建模目标说明，其中存在多处可能导致严重问题的错误或遗漏。请逐条指出问题所在，并说明每处问题可能造成的后果。

\begin{quote}
任务：预测下一小时各路段是否会发生拥堵。

数据：使用过去 6 个月的检测器数据，字段包括速度、流量和占有率。

方法：随机划分 80\% 训练集和 20\% 测试集。以准确率为评价指标，目标达到 90\% 以上。

特征：使用所有可用字段，包括当前时刻的速度、流量和占有率，以及这些字段在接下来 30 分钟的均值（通过提前计算）。

标签：速度低于 30 km/h 定义为拥堵（标签为 1），其余为不拥堵（标签为 0）。
\end{quote}

\subsection*{练习三：将现实问题翻译成建模定义}

选择以下两个情境之一，完成从模糊现实需求到完整建模定义的翻译过程。

\textbf{情境 A：缓解学校周边早高峰拥堵}

某城市交通管理部门希望通过数据分析，为学校周边的早高峰拥堵问题提供可操作的解决方案。可用数据：学校周边 15 个检测器的过去两年历史数据（5分钟粒度）；学校上课日历（开学、放假、考试日期）；周边 3 公里内的停车场占用记录；以及来自网约车平台的网格化出行需求数据。

\textbf{情境 B：提升地铁通勤体验}

某地铁公司希望通过数据分析降低早晚高峰的乘客投诉率。现有数据：过去一年的闸机进出站记录；列车实际到发时刻表（含延误信息）；乘客投诉记录（有文字描述，但无结构化分类）；以及每日运营报告（列车满载率、换乘等待时间）。

针对所选情境，完成以下工作：

\begin{enumerate}
  \item 识别至少两种不同的建模方向（例如，同一目标可以被翻译为预测问题或聚类问题），说明每种方向的适用条件和局限。
  \item 选择其中一种建模方向，完整填写五个维度（分析单元、任务类型、标签定义与来源、时空粒度、错误成本）。
  \item 列出至少三条必须在 Prompt 中明确声明的领域隐含知识。
  \item 将以上内容整合为一条完整的建模定义 Prompt，使 AI 能够在正确的问题框架下开展建模工作。
\end{enumerate}

\subsection*{练习四：评估代理指标的风险}

以下是三个常见的建模目标设定，每个都使用了代理指标代替真实目标。对每个目标，完成以下分析：（a）真实目标是什么，代理指标是什么；（b）在什么条件下优化代理指标会导致真实目标改善；（c）在什么条件下两者会出现分离，即代理指标提升但真实目标没有改善甚至恶化；（d）如果要补充一个验证机制来防止这种分离，你会选择什么。

\smallskip
\textbf{目标一：}通过预测模型优化出租车派单策略，以最小化“空驶率”（无乘客行驶里程占总行驶里程的比例）为模型优化目标，期望最终提升运营效率和乘客体验。

\smallskip
\textbf{目标二：}建立公交司机驾驶行为评分系统，以“急刹车次数”和“急加速次数”为主要指标，期望通过降低这两项指标来提升行车安全性和乘客舒适度。

\smallskip
\textbf{目标三：}建立高速公路巡逻路线优化模型，以“被巡逻路段比例”最大化为目标，期望提升路网安全管理覆盖率。

\subsection*{练习五：为AI设计完整的建模问题定义}

以下是一个真实研究场景的简要描述：

\begin{quote}
某研究团队拥有某城市 2020—2023 年的公交 GPS 数据（每 30 秒记录一次位置、速度和线路信息）、公交站点信息（位置坐标、所属线路、是否设有候车亭）、公交刷卡记录（上车时间、线路、车牌，无下车信息）以及同期的交通管理事件记录（施工、活动、特殊管制）。研究目标是：识别影响公交运营可靠性的主要因素，并据此提出改进建议。
\end{quote}

针对此场景，完成以下工作：

\begin{enumerate}
  \item 写出至少三种不同的建模方案，每种方案对应不同的任务类型，并说明每种方案可以回答的具体问题以及主要局限。
  \item 从三种方案中选择你认为最适合用于“提出改进建议”这一最终目的的方案，并为其完整设计建模定义（五个维度均需覆盖）。
  \item 写出一条完整的建模问题定义 Prompt，包括五要素（角色、上下文、任务、约束、输出格式），使 AI 能够直接依据该 Prompt 开展建模工作。
  \item 说明你在建模定义中做出的最重要的一个价值取舍（例如，你选择了某种标签定义而放弃了另一种），并解释为什么。
\end{enumerate}





